
%\documentclass{aa} 
\documentclass{aa} % Use the official aa class from CTAN
\usepackage{graphicx}
\usepackage{txfonts}
\usepackage{lipsum}
\usepackage{subcaption}     
\usepackage{lscape}       
\usepackage{placeins}   
\usepackage[round]{natbib}


\begin{document}

  \title{Gamma ray emission}
  \subtitle{Plasma conditions and physical processes}

  \author{Silvia Parente
    }

  \institute{University of Belgrade, MASS}

  \date{May 26th 2026}


 \abstract
 {
This review synthesizes the physical conditions and radiative processes governing gamma ray emission in extreme astrophysical environments, with a focus on active galactic nuclei and gamma ray bursts. The aim is to trace the evolutionary paths of high energy plasmas by examining non thermal continua, given by synchrotron and inverse compton scattering in relativistic collisionless shocks, and discrete thermal or nuclear signatures. The mechanisms of particle acceleration and the structural limitations of diffusive shock acceleration against the paradigm of magnetic reconnection in Poynting flux dominated outflows was evaluated and also the diagnostic utility of spectral architecture, using deviations from the Band function, nuclear spallation markers, and photoionized absorption features to map the chemical and thermodynamic state of the ejecta. Details about the modern multi-messenger observational era, the integration of space borne monitoring, ground-based Cherenkov arrays, and neutrino gravitational wave correlations are given to discuss bulk Lorentz factors, circumburst density profiles, and the physical limits of relativistic plasma dynamics.
}

\keywords{gamma rays --
        plasma --
        spectroscopy
        }

\maketitle
%%%%%%%%%%%%%%%%%%%%%%%%%%%%%%%%%%%%%%%%%%%%%%%%%%%%%%%%%%%%%%

\section{Introduction}

Gamma ray emissions represent the most energetic manifestation of electromagnetic radiation in the universe, and are useful to study environments characterized by extreme gravity, relativistic motion, and high temperature plasma dynamics. Phenomena such as gamma ray bursts (GRBs) and active galactic nuclei (AGNs) work as laboratories where energy releases can reach $10^{55}$ erg (isotropically equivalent) within seconds. This scale of energy density makes it necessary for relativistically expanding plasmas to exist, often collimated into narrow jets moving at Lorentz factors ($\Gamma$) of several hundred. Understanding the physical processes of these emissions requires a mix of plasma instability theory, nuclear reaction networks, and relativistic shock physics. 
The radiation field in these sources is usually divided into non thermal continua and discrete spectral features. Non thermal emission is mainly attributed to synchrotron radiation from electrons accelerated in relativistic collisionless shocks, that are mediated by collective plasma effects, such as the Weibel (filamentation) instability, that generates the turbulent magnetic fields necessary forhigh energy particle acceleration. In dense environments like optically thin accretion flows around black holes or subrelativistic supernova shocks, ion temperatures can exceed $10^{10}$ K: in hot plasmas like these, the emission is dominated by prompt nuclear deexcitation lines (0.1–10 MeV) and the decay of neutral pions ($\pi^0$) produced in inelastic hadronic collisions. 
Modern astrophysics has transitioned into a multimessenger era, where gamma ray data are complemented by detections of neutrinos, cosmic rays, and gravitational waves. This approach leads to the investigation of fundamental physics beyond the standard model, including searches for quantum gravity signatures like modified dispersion relations that may manifest over cosmological distances. However, interpreting these signals is still a challenge: spectral features originally thought to be cyclotron lines or annihilation peaks are now reevaluated in the context of photoionized plasma blobs entrained in relativistic fireballs. 
This review synthesizes current research on the physical conditions of gamma ray emitting plasmas; it explores the interplay between magnetic field generation, chemical evolution through spallation, and the radiative processes that shape observed spectra.

\section{Radiative Mechanisms}

\subsection{Non thermal continuum}
The main feature of GRB and AGN spectra is a broad non thermal continuum that spans over several orders of magnitude in energy. This emission originates from the dissipation of bulk kinetic energy within relativistic outflows. When a jet moving at a Lorentz factor $\Gamma \gg 1$ encounters either the ambient medium (external shocks) or faster moving shells within the jet itself (internal shocks), the resulting collisionless shocks become the main engines for particle acceleration. Particle acceleration in these environments is attributed, in general, to diffusive shock acceleration (DSA), or the first order Fermi process. As electrons cross and recross the shock front, they gain energy, eventually forming a non thermal power law distribution:$\frac{dN_e}{d\gamma_e} \propto \gamma_e^{-p} \quad \text{for} \quad \gamma_e > \gamma_m$ (where $\gamma_e$ is the electron Lorentz factor and $p$ typically ranges between 2.2 and 2.5). But the efficiency of this process hinges on the existence of a robust magnetic field. In the standard fireball model the magnetic field is often assumed to be a fraction $\epsilon_B$ of the shock’s internal energy. Modern plasma theory suggests these fields are generated locally via the Weibel (filamentation) instability, where counter streaming plasma beams create small scale magnetic turbulence that grows a lot on plasma timescales. Accelerated, these relativistic electrons spiral around the magnetic field lines, emitting synchrotron radiation; the characteristic frequency of this emission, $\nu_m$, corresponds to the minimum Lorentz factor $\gamma_m$ of the accelerated distribution. In the highly magnetized environments of GRB internal shocks the cooling timescale for these electrons is often shorter than the dynamical timescale of the expansion, and this leads to a "fast cooling" regime where the spectrum below the peak frequency follows a specific slope, $F_\nu \propto \nu^{-1/2}$, before transitioning to the harder $F_\nu \propto \nu^{-(p-1)/2}$ at higher energies. Synchrotron radiation provides the seed photons for the second major non thermal component: Inverse Compton (IC) scattering. In Synchrotron self compton (SSC) processes, the same population of relativistic electrons upscatters their own synchrotron photons to GeV and even TeV energies; the efficiency of this upscattering is described by the Compton $Y$ parameter, defined as the ratio of the IC power to the synchrotron power. While the Thomson cross section $\sigma_T$ governs scattering at lower energies, the process enters the Klein-Nishina (KN) regime when the energy of the photon in the electron's rest frame approaches the electron's mass energy ($m_e c^2$). In the KN regime, the scattering cross section drops:

\begin{equation}
    \sigma_{KN} \approx \sigma_T \frac{3}{8 \epsilon} \left( \ln(2\epsilon) + \frac{1}{2} \right)
\end{equation}

(where $\epsilon = h\nu / m_e c^2$). This reduction in efficiency results in a softening in thehigh energy gamma ray spectrum, providing a critical diagnostic for the energy density of the emitting region, and if the jet is embedded in a dense radiation field—such as photons originating from an accretion disk or a surrounding torus (External Inverse Compton, or EIC) the resulting spectrum is shifted and modified. The interplay between synchrotron, SSC, and EIC explains the two peaks in the spectral energy distribution (SED) frequently observed in blazars, where the first peak resides in the radio to optical range and the second peak dominates the gamma ray band. This non thermal continuum maps the history of energy transfer from the bulk motion of the plasma to the microscopic scales of individual leptons.

\subsection{Thermal and nuclear line emissions}

The non thermal continuum described in the previous section dominates the energy budget of relativistic jets, but the underlying plasma conditions in accretion disks and dense environments frequently give rise to different thermal and nuclear signatures. In these regions, especially near the event horizons of black holes or within sub relativistic shocks, the ion temperature $T_i$ can decouple from the electron temperature $T_e$. This decoupling allows $T_i$ to reach values higher than $10^{10}$ K, a regime where the kinetic energy of the ions is sufficient to overcome the Coulomb barrier of heavier nuclei (\citet{Kafexhiuetal2019}, p. 1). At these extreme temperatures inelastic collisions between protons, alpha particles and heavier nuclei like C and O trigger nuclear excitations. The main observational consequences are de excitation lines. Specifically, the 4.438 MeV line from $^{12}\text{C}^*$ and the 6.129 MeV line from $^{16}\text{O}^*$ are indicators of the plasma's chemical makeup (\citet{Kafexhiuetal2019}, p. 2). However, these lines don't exist in a vacuum, but they are subject to the cumulative effects of the plasma's nuclear history. As the plasma evolves, a process of nuclear spallation occurs where heavier nuclei are systematically broken down into lighter elements, protons and neutrons, which in turn alters the gamma ray emissivity in time (\citet{Kafexhiuetal2019}, p. 4). The presence of these features is complicated by the dynamic state of the plasma. In the context of GRBs, early spectral analysis identified narrow features in the 10–100 keV range, which were interpreted as cyclotron resonant scattering or $e^+-e^-$ annihilation lines (\citet{Haileyetal1999}, p. 1); while the cosmological shift and the relativistic nature of these sources have made these interpretations more complex, the entrainment of metal rich bubbles or filaments within the outflow remains a viable explanation for these features. These blobs, if photoionized by the intense radiation field of the fireball, can produce a range of absorption and emission features depending on the ionization parameter and the column density of the entrained material (\citet{Haileyetal1999}, p. 2).Furthermore, as the plasma temperature approaches the hadronic limit, the production of pions becomes a dominant cooling channel. The decay of neutral pions ($\pi^0 \rightarrow 2\gamma$) generates a broad hump in the spectrum centered around 70 MeV in the rest frame (\citet{Kafexhiuetal2019}, p. 1). Unlike the sharp de excitation lines, this hadronic signature is relatively insensitive to the specific chemical composition of the plasma, and instead depends on the proton density and the Maxwellian distribution of the ion energies. In pair dominated plasmas, which are common in thehigh energy environments of AGNs, the electron-positron annihilation line at $511$ keV (rest frame) provides a direct probe of the pair creation rate and the plasma's optical depth. In a moving jet, though, this line is typically Doppler boosted and broadened by both thermal motion and bulk expansion, causing it to blend into the power law continuum (\citet{Boettcheretal1996}, p. 2). The transition from a line dominated spectrum to a smooth continuum marks a shift in the physical conditions from a dense, thermalized accretion environment to a more diffuse, kinetically dominated relativistic flow.

\subsection{Hadronic decay channels}
When the ion temperature of the plasma approaches or exceeds $10^{11}$ K, or when relativistic protons are injected into a dense ambient medium, hadronic interactions become a primary channel forhigh energy gamma ray production (\citet{Kafexhiuetal2019}, p. 1). If leptonic processes are sensitive to low scale magnetic field perturbations, hadronic processes are governed by the strong force and require substantial threshold kinetic energies. The most prominent of these processes is the production and subsequent decay of neutral pions ($\pi^0$).
In a hot, thermal plasma dominated by protons, inelastic proton-proton ($pp$) collisions occur once the thermal kinetic energy overcomes the rest mass limit for pion production:
\begin{equation}
  p + p \rightarrow p + p + \pi^0
\end{equation}
The neutral pion is unstable and almost immediately ($\tau \approx 8.4 \times 10^{-17}$ s) decays into two photons:
\begin{equation}
  \pi^0 \rightarrow \gamma + \gamma
\end{equation}
Since the pion decays in flight, the resulting gamma ray spectrum is kinematically constrained; in the rest frame of the pion, each photon carries half of the pion's rest mass energy, $m_{\pi^0} c^2 / 2 \approx 67.5$ MeV. When integrated over a Maxwellian distribution of relativistic protons, this process produces a distinct, symmetric spectral hump centered at ~67.5 MeV when plotted on a linear scale, or a broad feature spanning from tens of MeV to several GeV (\citet{Kafexhiuetal2019}, p. 1, 6). The shape and amplitude of this hadronic feature depend on the underlying proton temperature and density more than the electron distribution, making it an unambiguous indicator of baryonic acceleration (\citet{Kafexhiuetal2019}, p. 7).

In more complex relativistic environments, like the highly collimated jets of blazars or the internal shock zones of gamma ray bursts, protons can also interact with ambient photon fields. This process, known as photo-pion production or the $p\gamma$ channel, becomes dominant when the product of the proton energy and the photon energy exceeds the $\Delta(1232)$ resonance limit. The primary resonant channels are:

\begin{equation}
  p + \gamma \rightarrow \Delta^+ \rightarrow p + \pi^0
\end{equation}
\begin{equation}
  p + \gamma \rightarrow \Delta^+ \rightarrow n + \pi^+
\end{equation}

The neutral pions generated via this resonant state decay into gamma rays while the charged pions ($\pi^+$) decay into positrons and neutrinos 

\begin{equation}
    (\pi^+ \rightarrow \mu^+ + \nu_\mu \rightarrow e^+ + \nu_e + \nu_\mu + \bar{\nu}_\mu).
\end{equation}

The transition from $pp$ dominated emission to $p\gamma$ dominated emission depends heavily on the radiation field density within the plasma zone. In dense accretion flows, the high column density of matter favors thermal $pp$ collisions (\citet{Kafexhiuetal2019}, p. 1). In the diffuse, radiation dominated regions of relativistic fireballs, the massive flux of synchrotron or accretion disk photons makes $p\gamma$ resonances the primary mechanism for breaking downhigh energy proton streams (\citet{Boettcheretal1996}, p. 1); the gamma rays produced via these hadronic channels do not always escape freely, but they frequently interact with the same target photon fields that produced them, initiating electron-positron pair cascades that redistribute thehigh energy budget down to the leptonic continuum (\citet{Boettcheretal1996}, p. 2; \citet{KumarandZhang2015}, p. 29).

\subsection{Pair plasma dynamics and gamma ray photosphere}
The high luminosity ofhigh energy transients, when coupled with their rapid variability timescales, imposes that the primary emission originates from extremely compact regions. If these regions were stationary, the sheer density ofhigh energy photons would result in catastrophic photon-photon collision rates; when the center of mass energy of two interacting photons exceeds the limit of twice the electron rest mass, they undergo pair production via 
\begin{equation}
    \gamma + \gamma \rightarrow e^+ + e^-
\end{equation}.
This compactness problem makes the emitting plasma undergo extreme relativistic expansion; otherwise, the source would be completely opaque, suppressing thehigh energy flux we observe entirely (\citet{KumarandZhang2015}, p. 6).
In the innermost regions of an active galactic nucleus jet or the early stages of a gamma ray burst fireball, the creation of these electron-positron pairs dominates the local plasma dynamics. The lepton density drives the Thomson optical depth to ($\tau \gg 1$) and traps the radiation within the flow. The plasma only becomes transparent to its own emission at the photospheric radius, which is mathematically the boundary where the scattering optical depth drops to $\tau \approx 1$(\citet{Peer2019}, p. 9; \citet{KumarandZhang2015}, p. 15); here the trapped radiation decouples from the baryons and the pair plasma. Because the photons and leptons were previously in strict thermal contact through rapid compton scattering, the escaping radiation shows a quasi thermal signature. Recent spectral analyses of prompt GRB emissions frequently reveal this component that appears as a modified blackbody spectrum superimposed on the standard non thermal power law continuum (\citet{Peer2019}, p. 10). The peak energy of this photospheric emission is determined by the initial temperature at the base of the flow, blueshifted by the jet's bulk Lorentz factor $\Gamma$.
However, a pure blackbody profile is rarely observed in isolation. If bulk kinetic energy or magnetic energy is dissipated deep within the optically thick regime, for instance through deeply embedded internal shocks or localized magnetic reconnection events, the plasma is reheated: the radiation field interacts with any newly accelerated electrons through inverse Compton scattering and this process, known as sub photospheric dissipation, broadens the thermal peak and hardens thehigh energy tail of the spectrum, blurring the boundary between thermal release and non thermal particle acceleration (\citet{Peer2019}, p. 11).
As the outflow expands beyond the photosphere and the temperature drops below the relativistic limit, the pair creation rate drops. The evolutionary trajectory of the plasma then becomes governed by the balance between adiabatic expansion cooling and pair annihilation kinetics; in the confined jets of AGNs the continuous evolution of the $e^- - e^+$ distribution is dictated by radiative energy losses and the cross section of annihilation (\citet{Boettcheretal1996}, p. 1-2). One might expect localized annihilation features (like a 511 keV peak in the comoving frame) to emerge as the pairs recombine, the relativistic bulk motion and severe thermal broadening typically smear these discrete lines into the broader continuum. Because of that direct detection of annihilation lines remains an observational challenge, acting as generalized absorption or emission features depending on the immediate column density of the entrained plasma (\citet{Boettcheretal1996},\citet{Haileyetal1999}, p. 1).

\section{Observational framework}
\subsection{Space borne detectors}
The understanding of gamma ray emitting plasmas is fundamentally tied to the capabilities of space borne instrumentation: since the Earth's atmosphere is opaque tohigh energy photons, tracking the evolution of transients like GRBs and AGNs requires satellite based platforms. 
Early orbital observations, particularly by instruments like the Burst and Transient Source Experiment (BATSE) aboard the Compton Gamma Ray Observatory, established that prompt gamma ray emission occurs on timescales of milliseconds to hundreds of seconds (\citet{KumarandZhang2015}, p. 2). BATSE’s field of view was optimized for detection and localization, mapping out an isotropic distribution that hinted at a cosmological origin; however, thehigh energy plasma states during this initial phase are temporary. The rapid spectral variability observed by these early instruments revealed a non thermal 'band' function peak, which modern theory interprets as either synchrotron emission from internal shocks or modified thermal radiation emerging from a dynamic, pair dominated photosphere (\citet{Peer2019}, p. 10; \citet{KumarandZhang2015}, p. 15).
A major progress occurred with the launch of the Swift satellite (later renamed the 'Neil Gehrels Swift Observatory'). Swift combined a wide field coded aperture mask, the Burst Alert Telescope (BAT), with two narrow field instruments: the X-ray Telescope (XRT) and the Ultraviolet/Optical Telescope (UVOT)(\citet{KumarandZhang2015}, p. 2). This capability of the spacecraft allowed it to reposition its narrow field instruments toward a burst within less than a minute of the initial prompt trigger, and this response closed the observational gap between the energetic flash and the early afterglow, revealing features in the light curves like steep early decay phases, shallow plateau phases, and intense X-ray flares (\citet{KumarandZhang2015}, p. 30).
These space borne observations show that the transition from prompt emission to the afterglow marks a fundamental shift in the place of plasma dissipation: while the prompt phase is dominated by internal physics, such as shell collisions or magnetic reconnection events within the relativistic outflow, the afterglow is initiated when the jet plows into the surrounding medium (\citet{Peer2019}, p. 1; \citep{Brainerd1999} and as the expanding shell sweeps up external matter, it drives a forward shock into the ambient gas and a reverse shock back into the ejecta. 
Space borne monitoring via Swift and the Fermi gamma ray Space Telescope has shown that this interaction causes the bulk Lorentz factor $\Gamma$ to decay as a power law over time:
$\Gamma(t) \propto t^{-3/8}$
This self similar deceleration blast wave, described by the Blandford McKee solution, causes the characteristic synchrotron frequencies to sweep downward through the X-ray, ultraviolet, optical, and radio bands over days and weeks (\citet{KumarandZhang2015}, p. 32).

Furthermore, Fermi's Large Area Telescope (LAT) extended detection capabilities into the GeV range, revealing thathigh energy gamma ray emission often goes far longer than the MeV prompt phase (\citet{KumarandZhang2015}, p. 27). The GeV component is interpreted ashigh energy inverse Compton scattering occurring at the forward shock front, where electrons continue to be accelerated by collective plasma instabilities \citep{Brainerd1999}. 

\subsection{Ground based Cherenkov arrays}
As the photon energy budget ofhigh energy cosmic accelerators extends past the GeV scale and enters the TeV to PeV regimes, space borne tracking becomes geometrically unfeasible. Because of the power law decline of particle distributions the photon flux at these extreme energies drops so low that the collecting areas of satellite detectors, typically under $1\text{ m}^2$, cannot secure statistically significant counts (\citet{KumarandZhang2015}, p. 2). Resolving these ultrahigh energy plasma states requires using the Earth's atmosphere as a natural calorimeter: ground based detection methodologies leverage the interactions of thesehigh energy photons with atmospheric nuclei, a process that starts an extensive air shower (EAS).
When a primary gamma ray photon enters the upper atmosphere, it interacts with the electrostatic field of an atmospheric nucleus and undergoes pair production ($\gamma \rightarrow e^+ + e^-$). The relativistic electron and positron then radiatehigh energy secondary photons via bremsstrahlung as they go close to other nuclei. This cycle of pair production and bremsstrahlung repeats rapidly and generates a cascading avalanche of millions of secondary leptons and photons traveling downward at near the speed of light (\citet{Addazietal2021}, p. 23). Since these secondary particles move through the atmosphere faster than the local phase velocity of light ($c/n$, where $n$ is the refractive index of air) they polarize the surrounding gas molecules, which then emit a collimated beam of blue to ultraviolet Cherenkov radiation.
Imaging Atmospheric Cherenkov Telescopes (IACTs), such as H.E.S.S., MAGIC, and VERITAS, are designed to capture this fleeting flash of Cherenkov light, which lasts just a few nanoseconds. These systems can reconstruct the cascade's geometry from multiple viewing angles by deploying arrays of large optical reflectors that have ultra fast photomultiplier tube cameras; this approach allows researchers to determine the arrival direction of the primary gamma ray and estimate its energy based on the total light yield. More importantly, it provides a means to differentiate between gamma ray induced showers and the much more numerous showers triggered by hadronic cosmic rays (\citet{Addazietal2021}, p. 23; \citet{Boettcheretal1996}, p. 1): hadronic showers, mediated by the strong interaction, are characterized by large transverse momenta and a high fraction of pions, resulting in a chaotic and asymmetric shower structure while purely electromagnetic showers are highly compact and symmetric instead.
Ground based experimentalists also utilize Extensive Air Shower arrays, such as HAWC or LHAASO: instead of focusing on the optical Cherenkov light from the air column, these arrays place thousands of water Cherenkov detectors directly on the ground to sample the secondary particles that survive the trip to the surface (\citet{Addazietal2021}, p. 23) and when the secondary leptons plow through these pure water tanks, they produce localized Cherenkov flashes that are recorded by submerged sensors.
The data gathered through these Ground based techniques have changed our understanding of the radiative processes occurring within the relativistic jets of active galactic nuclei. Observations of rapid TeV variability in blazars have forced a reassessment of the spatial scales of emission zones and the role of the Klein-Nishina suppression effect inhigh energy inverse Compton scattering (\citet{Boettcheretal1996}, p. 1-2). Since TeV photons are highly susceptible to pair production when interacting with lower energy photons, the very detection of these ground based signals acts as a limit on the opacity of the source's local plasma environment and the diffuse extragalactic background light (EBL) (\citet{Boettcheretal1996}, p. 2; \citet{KumarandZhang2015}, p. 6).

\subsection{Multimessenger point of view}
High energy photons are sensible to localized absorption, pair production, and scattering, that can obscure the innermost engines of AGNs and GRBs (\citet{Boettcheretal1996}, p. 2; \citet{KumarandZhang2015}, p. 6). A multi messenger framework, that integrates gamma ray spectroscopy with neutrino detections and gravitational wave registration, provides an unobstructed window into the densest, most highly magnetized regions of these cosmic accelerators (\citet{Addazietal2021}, p. 24).
High energy astrophysical neutrinos are a signature of hadronic acceleration within the plasma: while leptons produce the observed continuum via synchrotron and inverse Compton, they can't generate cosmic neutrinos. As established in Section 2.3, when a relativistic jet contains a huge baryonic component, accelerated protons collide with ambient photons or target gas (\citet{Boettcheretal1996},\citet{Kafexhiuetal2019}, p. 1) and these photo-pion ($p\gamma$) and proton-proton ($pp$) interactions yield charged pions ($\pi^\pm$) alongside neutral pions (\citet{KumarandZhang2015}, p. 28). The charged pions subsequently decay through the leptonic channel:
\begin{equation}
    \pi^+ \rightarrow \mu^+ + \nu_\mu \rightarrow e^+ + \nu_e + \nu_\mu + \bar{\nu}_\mu
\end{equation}

Neutrinos interact just through the weak force, so they escape the dense plasma architecture of the jet or accretion flow without attenuation. Detecting a TeV–PeV neutrino in spatial and temporal coincidence with a gamma ray flare confirms that the underlying plasma is acting as a cosmic ray factory, pinning down the baryonic energy fraction of the blast wave (\citet{KumarandZhang2015}, p. 28; \citet{Addazietal2021}, p. 24).

Simultaneously, gravitational waves provide a completely non electromagnetic metric for the dynamical initiation of these plasmas, because in short duration GRBs the merger of binary neutron stars or a neutron star-black hole system distorts spacetime, radiating gravitational waves that map the system's mass and orbital evolution prior to fusion. It leaves behind a hypermassive neutron star or a stellar mass black hole surrounded by a dense, highly heated accretion torus: this extreme environment generates the high temperature pair plasmas and collimated, magnetically dominated outflows discussed in the context of photospheric emission (\citet{Peer2019}, p. 10; \citet{KumarandZhang2015}, p. 15).

Correlating the arrival time of the gravitational wave chirp with the first registration of gamma rays places limits on fundamental physics. For example, the joint detection of GW170817 and GRB 170817A established that the speed of gravity matches the speed of light to within one part in $10^{15}$; this temporal correlation constrains the jet propagation time through the small amount of material ejected during the merger's early tidal disruption phase. On larger scales, these correlated signals allow physicists to test quantum gravity alternatives, such as Lorentz invariance violation (LIV). If the smooth structure of spacetime breaks down at the Planck scale,high energy gamma rays would experience an energy dependent velocity shift over cosmological distances (\citet{Addazietal2021}, p. 12, 23). Researchers can place bounds on modified dispersion relations by comparing the temporal features of gamma ray light curves against the absolute arrival time of gravitational waves or low energy neutrinos:

$E^2 = p^2 c^2 \left[ 1 \pm \left( \frac{E}{E_{\text{QG}}} \right)^\alpha \right]$

where $E_{\text{QG}}$ is the quantum gravity energy scale and $\alpha$ dictates the order of the correction (\citet{Addazietal2021}, p. 11, 23). 

\section{Spectral architecture}
\subsection{The Band function}

The spectral energy distributions of GRBs have historically been modeled using the band function. This model joins two distinct power laws: ahigh energy photon index $\alpha$ and ahigh energy photon index $\beta$, transitioning at a characteristic break energy denoted as $E_0$ or peak energy $E_p$(\citet{KumarandZhang2015}, p. 11). The functional form is expressed as:
\begin{equation}
    N_E(E) = A \left( \frac{E}{100 \text{ keV}} \right)^\alpha \exp\left(-\frac{E}{E_0}\right) \quad \text{for} \quad E \le (\alpha - \beta)E_0
\end{equation}
\begin{equation}
    N_E(E) = A \left[ \frac{(\alpha - \beta)E_0}{100 \text{ keV}} \right]^{\alpha - \beta} \exp(\beta - \alpha) \left( \frac{E}{100 \text{ keV}} \right)^\beta \quad \text{for} \quad E > (\alpha - \beta)E_0
\end{equation}
Statistical compilations of many bursts reveal that thehigh energy index converged around $\alpha \approx -1$, while thehigh energy index typically settled near $\beta \approx -2.25$ (\citet{KumarandZhang2015}, p. 11). The band function provides a good fit to broad band data from instruments like BATSE and Fermi but it lacks a physical basis and this complicates efforts to identify the plasma conditions responsible for the emission.

The main physical conflict arises when comparing the observedhigh energy index $\alpha$ against the theoretical limits of synchrotron radiation: standard analytical treatments of particle acceleration in collisionless shocks say that if the accelerated electrons are cooling rapidly, as expected in highly magnetized internal shocks, thehigh energy spectrum can't be harder than $F_\nu \propto \nu^{-1/2}$, corresponding to a photon index $\alpha = -1.5$(\citet{Peer2019}, p. 9). Even in the cooling regime, where the lower boundary is set by the pitch angle averaged synchrotron spectrum of a single electron, the maximum hardness is limited to $F_\nu \propto \nu^{1/3}$, or $\alpha = -0.67$ (\citet{Peer2019}, p. 9; \citet{KumarandZhang2015}, p. 12). A big fraction of observed GRB spectra exhibit $\alpha$ values warmer than $-1.5$, violating the fast cooling prediction.
To reconcile these deviations, models must cgange the simple power law assumption. High resolution spectroscopic data from the Fermi gamma ray Space Telescope have revealed deviations from the classical band function, a localized, sub dominant thermal component or an additionalhigh energy power law (\citet{KumarandZhang2015}, p. 14, 27); when a blackbody component with a temperature $k_B T \approx 10\text{–}100\text{ keV}$ is integrated into the spectral fit, the apparent value of $\alpha$ softens and pulls the underlying non thermal index below the synchrotron line of death. This joint modeling suggests that the prompt spectrum is a composite structure that joins non thermal shock emission with a hot, quasi thermal photospheric component escaping from the interior of the fireball (\citet{Peer2019}, p. 10; \citet{KumarandZhang2015}, p. 15).
Systematic deviations at the highest energies often present as an extra power law component extending in the GeV: this unabsorbedhigh energy tail is a clue of secondary physical processes, like early time inverse Compton upscattering or hadronic cascades taking place within the diffuse outer zones of the expanding jet (\citet{Boettcheretal1996}, p. 1; \citep{Brainerd1999}. In ultra compact sources thehigh energy spectrum can display a steepening or cutoff that deviates below the predicted $\beta$ index: this is a direct diagnostic of localized pair production opacity ($\gamma\gamma \rightarrow e^+e^-$), and allows physicists to map the physical dimensions and bulk Lorentz factor of the emitting plasma blob (\citet{KumarandZhang2015}, p. 6; \citet{Boettcheretal1996},\citet{Haileyetal1999}, p. 1).

\subsection{Nuclear deexcitation lines, spallation and chemical evolution}
High temperature baryonic components in dense plasma regions produce a distinct forest of nuclear de excitation lines (\citet{Boettcheretal1996},\citet{Kafexhiuetal2019}, p. 1). This typically happens in the inner regions of accretion disks feeding supermassive black holes or inside sub relativistic shocks, where the ion temperature $T_i$ climbs above $10^{10}\text{ K}$ (\citet{Boettcheretal1996},\citet{Kafexhiuetal2019}, p. 1) and where the kinetic energy distributed among the ions is sufficient to overcome the Coulomb barriers of ambient metals, resulting in inelastic collisions that excite nuclei.

The first consequence of thesehigh energy collisions is the isotropic emission of narrow gamma ray lines as the excited states decay to ground configurations, like the $4.438\text{ MeV}$ line from $^{12}\text{C}^*$ and the $6.129\text{ MeV}$ line from $^{16}\text{O}^*$ (\citet{Boettcheretal1996},\citet{Kafexhiuetal2019}, p. 2). Additional spectral tracking can identify lines from heavier elements, like the $0.847\text{ MeV}$ line of $^{56}\text{Fe}^*$, the $1.635\text{ MeV}$ line of $^{20}\text{Ne}^*$, and the $2.313\text{ MeV}$ line of $^{14}\text{N}^*$. These lines are localized in energy, so their profiles acts as a thermometer for the plasma’s ion temperature, while their cumulative flux measures the relative abundances (\citet{Boettcheretal1996},\citet{Kafexhiuetal2019}, p. 2).
However, these abundance profiles are not static: the samehigh energy collisions that induce nuclear excitation also induce a continuous process of chemical destruction through nuclear spallation (\citet{Boettcheretal1996},\citet{Kafexhiuetal2019}, p. 4). When a high temperature proton or alpha particle strikes a heavy nucleus like carbon, oxygen or iron, the interaction breaks the target nucleus apart rather than merely exciting it and this depletes the initial reserve of heavier metals, transforming them into lighter elements:
\begin{equation}
    p + {}^{12}\text{C} \rightarrow p + p + {}^{11}\text{B}
\end{equation}
\begin{equation}
    p + {}^{16}\text{O} \rightarrow p + \alpha + {}^{12}\text{C}
\end{equation}

Over the lifetime of the plasma this spallation driven destruction creates a chemical evolution that can be monitored spectroscopically. As the primary targets like $^{12}\text{C}$ and $^{16}\text{O}$ are burned away, the gamma ray line emissivity profiles shift over time, the initialhigh energy metal lines fades while secondary lines from spallation products like Boron, Beryllium, and Lithium become more pronounced (\citet{Boettcheretal1996},\citet{Kafexhiuetal2019}, p. 4). This chemical evolution means that the observed ratio of the $4.438\text{ MeV}$ line to secondary lines allows as a tool to measure integrated exposure time of the hot plasma phase.
In the context of relativistic fireballs, the preservation or destruction of these chemical markers is dependent on the spatial stratification of the outflow: if metal rich plasma blobs or filaments get into the base of a GRB jet from the progenitor star, they are exposed to a brutal radiation field and intense internal collisions (\citet{Boettcheretal1996},\citet{Haileyetal1999}, p. 1). If the plasma remains compact and optically thick, the nuclear signatures are entirely washed out by compton scattering or undergo complete spallation down to raw protons and neutrons. However, if these photoionized filaments manage to survive or are injected into the lower density outer regions of the expansion, they can manifest as broad absorption or emission features in the X-ray and soft gamma ray bands (\citet{Boettcheretal1996},\citet{Haileyetal1999}, p. 2). 

\subsection{Compton attenuation and absorption features in photoionized plasmas} 
High energy photons propagating through a dense outflow rarely escape into the surrounding medium without interacting with the local environment. The surrounding plasma acts as a dynamic, highly opaque screen that continuously modifies the intrinsic spectrum described by the Thomson optical depth, governed by the density of ambient leptons. However, for energetic photons scattering off thermalized electrons, or lower energy seed photons being upscattered by a relativistic particle distribution, comptonization modifies the emergent radiation profile (\citet{Peer2019}, p. 9; \citet{Boettcheretal1996}, p. 1). When the photon energy in the rest frame of the electron approaches or exceeds $m_e c^2$, the scattering process enters the Klein-Nishina regime where the scattering cross section diminishes logarithmically with energy, that sharply reduces the cooling efficiency of the highest energy electrons and induces a noticeable steepening or attenuation in the high energy tail of the observed spectrum (\citet{Boettcheretal1996}, p. 1-2).
Localized plasma conditions can generate highly specific, discrete spectral structures. Early observational records of gamma ray bursts in the 10–100 keV band frequently exhibited narrow absorption and emission anomalies; these structures were initially categorized as cyclotron resonant scattering or electron-positron annihilation lines, because of the assumption that bursts originated from galactic neutron stars (\citet{Haileyetal1999}, p. 1). Once redshift measurements confirmed the cosmological distances of these transients, the energetic requirements rendered the local neutron star framework obsolete. The observed features, heavily blueshifted by the bulk Lorentz factor $\Gamma$ of the jet, highlighted the need for an alternative physical explanation.
The current theoretical model is that that these spectral anomalies originate from metal rich plasma blobs or dense filaments in the relativistic outflow: as the big radiation flux from the central engine overtakes these slower, denser clumps, the material is violently photoionized (\citet{Haileyetal1999}, p. 1). The thermodynamic state and opacity of this material are dependent on the ionization parameter, $\xi = L / n r^2$, where $L$ is the ionizing luminosity, $n$ is the local electron density, and $r$ is the radial distance from the emission source. If $\xi$ is exceptionally high, the heavier elements are fully stripped of their bound electrons, so that the plasma is completely transparent to atomic line absorption, but if the parameter is too low, boun-free opacity dominates and suppresses the soft X-ray and gamma ray continuum.
When the ionization parameter and the line of sight column density ($N_H$) are in the viable range, the photoionized filaments function as a spectral filter because the continuum radiation passing through these blobs is subjected to bound-free absorption and resonant line scattering, generating broadband structures that can heavily modify the spectral slope. These interactions manifest as K-shell absorption edges from highly ionized species, such as hydrogen like or helium like iron, and blended absorption troughs that emulate the appearance of single, broad lines in low resolution data (\citet{Haileyetal1999}, p. 2). Analyzing these compton attenuated spectra and photoionized absorption features works as a tool to measure the metallicity, column density, and spatial stratification of the matter surrounding the central engine.

\subsection{Propagation effects}
When gamma rays escape the local environment of the central engine, in their transit in intergalactic space, they get through erosion and temporal shifting. The diffuse bath of photons generated by stellar emission and dust reprocessing, EBL, acts as a scattering target: for a primary gamma ray with energy $E_\gamma$, a collision with a background photon of energy $\epsilon$ triggers pair production ($\gamma + \gamma \rightarrow e^+ + e^-$) if the center-of-mass energy exceeds the kinematic threshold:
$$s \approx 4 E_\gamma \epsilon \ge 4 m_e^2 c^4$$
Because the EBL density evolves with redshift, the resulting optical depth is a strong function of the source's distance and the observed photon energy both (\citet{Boettcheretal1996}, p. 2), so that the high energy spectra of distant blazars and GRBs are artificially steepened by the time they reach Earth. Studying and correcting for this redshift dependent attenuation is essential for recovering the intrinsic emission profile of the source plasma, establishing an energy dependent horizon beyond which the universe becomes opaque to TeV emission (\citet{KumarandZhang2015}, p. 6).
Propagation over billions of light years allows astrophysicists to test the structural integrity of spacetime itself: certain formulations of quantum gravity propose that at the Planck scale, spacetime ceases to be a continuous, smooth manifold and instead it exhibits a discrete, fluctuating geometry. This extreme quantization could lead to a violation of Lorentz invariance (LIV), causing the cosmic vacuum to behave as a dispersive medium for ultra high energy photons. Under LIV, a photon's group velocity becomes slightly dependent on its energy, altering the standard special relativistic dispersion relation to:
$$E^2 = p^2 c^2 \left[ 1 \pm \left( \frac{E}{E_{\text{QG}}} \right)^n \right]$$
where $E_{\text{QG}}$ represents the energy scale of quantum gravity (often theoretically anchored near the Planck energy, $\approx 1.22 \times 10^{19} \text{ GeV}$) and the exponent $n$ dictates the order of the quantum correction (\citet{Addazietal2021}, p. 11).
If LIV holds, photons of different energies emitted simultaneously from a distant source will gradually drift apart in time. Transients with high amplitude, rapid variability and cosmological redshifts are the optimal laboratory for isolating this effect: by cross correlating the arrival times of MeV photons with those of GeV or TeV photons originating from the same localized flaring event, researchers can search for these tiny, cumulative temporal offsets (\citet{Addazietal2021}, p. 23).
Isolating this vacuum dispersion is very challenging, because the intrinsic emission delays generated by particle acceleration timescales and shock physics in the source plasma must be precisely modeled and subtracted. Nevertheless, continuous monitoring of these delays places more stringent lower bounds on $E_{\text{QG}}$, restricting the parameter space for string theory and loop quantum gravity models (\citet{Addazietal2021}, p. 23). 

\section{Plasma diagnostics}

\subsection{Bulk Lorentz factors, jet opening angles}
The extreme physical environments characterizing GRBs and AGNs can't be modeled under the assumption of Newtonian expansion: the most fundamental parameter governing the kinematics of these sources is the bulk Lorentz factor, $\Gamma$, which dictates both the apparent luminosity and the timescale of the observed emission. 
During the prompt phase of a GRB, the central engine releases a massive flux of high energy photons over a variability timescale, $\delta t$, that can be as short as milliseconds; if the emission region were stationary or expanding non relativistically, the physical radius of the source would be restricted to $R \le c \delta t$. Because of the immense isotropic equivalent luminosity ($L_{iso}$), the photon density within this compact volume would be too high. Because the center of mass energy of these colliding photons would frequently exceed the threshold for pair production ($\gamma\gamma \rightarrow e^+ e^-$), the resulting optical depth to pair creation ($\tau_{\gamma\gamma}$) would be $10^{13}$ to $10^{15}$ times greater than unity (\citet{Peer2019}, p. 9) and under such conditions, no high energy gamma rays could escape. To reconcile this theoretical opacity with the actual detection of MeV and GeV photons, the plasma must be expanding relativistically toward the observer. This bulk motion scales the emission radius up by a factor of $\Gamma^2$ and shifts the comoving photon energies downward by a factor of $\Gamma$, severely reducing the available phase space for pair production. Imposing that $\tau_{\gamma\gamma} < 1$ establishes a theoretical lower limit for the outflow, requiring $\Gamma \gtrsim 100$ for GRB fireballs (\citet{KumarandZhang2015}, p. 6; \citet{Peer2019}, p. 9).
Direct estimates of the initial bulk Lorentz factor $\Gamma_0$ are obtained by monitoring the early onset of the afterglow. As the relativistic ejecta sweeps up mass from the circumburst medium, it drives a forward shock into the ambient gas. The blast wave continues to coast at approximately $\Gamma_0$ until the swept up mass, $M_{sw}$, divided by $\Gamma_0$, equals the initial rest mass of the ejecta; at this critical deceleration radius, the flow begins to shed kinetic energy rapidly. This hydrodynamic transition manifests observably as a smooth, distinct peak in the early optical or X-ray light curve at a specific deceleration time, $t_{dec}$ (\citet{KumarandZhang2015}, p. 30) and by measuring $t_{dec}$ alongside the total released energy and the estimated density of the ambient medium, astrophysicists can calculate the precise value of $\Gamma_0$ independently of opacity limitations.
Determining $\Gamma$ is only a part of the kinematic puzzle: identifying the true energy scale of the explosion requires measuring the jet opening angle, $\theta_j$ and if high energy transients exploded spherically, the energy budget required to power them would exceed the rest mass of a solar type star. However, these outflows are highly collimated; a relativistic jet focuses its radiation into a narrow beaming cone with a half angle of roughly $1/\Gamma$ due to relativistic aberration (\citet{KumarandZhang2015}, p. 32). As long as the flow remains ultra relativistic, this beaming cone is narrower than the physical boundaries of the jet ($1/\Gamma < \theta_j$), and the observer cannot distinguish the collimated outflow from a spherical explosion.
As the forward shock sweeps up matter and decelerates, $\Gamma$ decreases until eventually the relativistic beaming angle widens enough to encompass the physical edge of the jet ($1/\Gamma \approx \theta_j$). Now, the observer can detect a deficit in emission because there is no plasma beyond $\theta_j$ to contribute to the broadening beam and concurrently the loss of lateral confinement allows the jet to expand sideways, accelerating the dissipation of internal energy. This geometric transition produces a characteristic jet break, an achromatic steepening in the decay slope of the afterglow light curve across all electromagnetic bands (\citet{KumarandZhang2015}, p. 32). Identifying the exact timing of this break provides a direct mathematical pathway to $\theta_j$ and allows researchers to correct the artificially inflated $L_{iso}$ down to the corrected energy of the explosion, confirming that the central engines of GRBs release energies broadly clustered around a standard value of $10^{51}$ ergs.

\subsection{Magnetic field generation}
The presence of non thermal synchrotron radiation in the spectra of GRBs and blazars imposes this physical requirement: the emitting plasma must channel robust magnetic fields. In the classical fireball paradigm, the energy density of this field is parameterized as a fraction $\epsilon_B$ of the total internal energy directly behind a shock front. Since the pre shock interstellar medium or stellar wind has magnetic fields far too weak to account for the observed multi wavelength emission, the necessary field amplification must happen during the shock crossing itself (\citet{Peer2019}, p. 1).
In collisionless relativistic shocks, this rapid amplification is attributed to kinetic plasma phenomena: when a rapidly moving plasma shell overtakes the ambient medium or a slower internal shell, the lack of immediate collisionality causes the counterstreaming beams of charged particles to interpenetrate. This high velocity anisotropy triggers the filamentation, or Weibel, instability \citep{Brainerd1999} that swiftly channels a fraction of the directed bulk kinetic energy into localized, intense magnetic turbulence. This mechanism provides the chaotic magnetic scattering centers necessary to trap electrons near the shock boundary by operating on the microscopic scale of the plasma skin depth, $c/\omega_p$ (where $\omega_p$ is the plasma frequency). This localized trapping is the foundational requirement for first-order Fermi acceleration, allowing particles to cross the shock repeatedly and gain energy.
Despite its theoretical utility in initiating particle acceleration, the Weibel instability has limitations when applied to the full evolutionary track of the plasma: the microscopic fields generated by kinetic instabilities are transient and numerical simulations repeatedly indicate that these turbulence driven fields decay rapidly in the downstream region of the shock. If the magnetic field dissipates faster than the time required for electrons to cool via synchrotron emission, the shock fails to produce the non thermal continua observed in extended afterglows.
This theoretical tension has caused a major shift toward magnetohydrodynamic (MHD) formulations, moving away from purely thermally driven fireballs. If the central engine, whether a black hole accretion system or a highly magnetized magnetar, launches the jet with a dominant magnetic field, the outflow's dynamics are dictated by the Poynting flux rather than thermodynamic pressure (\citet{KumarandZhang2015}, p. 1): in these models, the initial state of the plasma is characterized by a high magnetization parameter, $\sigma$, defined as the ratio of magnetic energy density to the rest mass energy density of the outflow. As the jet expands, it does not rely primarily on internal hydrodynamic collisions to dissipate kinetic energy but instead it dissipates its stored magnetic energy directly.
The main engine for particle acceleration in these highly magnetized flows is magnetic reconnection (\citet{Peer2019}, p. 1). As the relativistic jet propagates outward, complex topological structures, magnetohydrodynamic instabilities, or turbulent motions force magnetic field lines of opposite polarity to interact. The structural collapse and reconfiguration of the magnetic field transfer magnetic energy to the entrained leptons. Magnetic reconnection populates the high energy tail of the electron distribution required to match the hardest prompt GRB observations. 

\subsection{Particle acceleration efficiency, diffusion limits}

In the standard framework of GRBs and AGNs, the observed non thermal continua needs a highly efficient mechanism for transferring macroscopic bulk kinetic energy into the microscopic kinetic energy of individual particles. The maib candidate for this energy transfer is diffusive shock acceleration (DSA), operating via the first order Fermi process. This means that charged particles are energized as they repeatedly cross and recross a relativistic shock front. Particles in relativistic collisionless shocks gain energy by a factor of roughly two with each complete back and forth crossing (\citet{KumarandZhang2015}, p. 20), unlike non relativistic shocks, where energy gains are marginal per cycle. However, the efficiency of this mechanism is constrained by the plasma's ability to magnetically confine the particles within the immediate vicinity of the shock.
The acceleration timescale is limited by the spatial diffusion of the particles through the upstream and downstream magnetic turbulence (\citet{Peer2019}, p. 4): in the most optimistic theoretical scenario, the Bohm diffusion limit, the scattering mean free path of a particle is reduced to its Larmor radius and under this extreme condition the minimum comoving time required to accelerate a particle of mass $m$ to a Lorentz factor $\gamma$ is equivalent to its Larmor time:
$$t'_{acc} \approx \frac{\gamma m c}{q B'}$$
where $B'$ is the comoving magnetic field strength (\citet{KumarandZhang2015}, p. 20).

While the Bohm limit sets the maximum theoretical rate of energy gain, the particles are subjected to severe radiative energy losses. For leptons crossing a highly magnetized plasma, synchrotron cooling acts as a brake on the acceleration process. The comoving cooling timescale for an electron is inversely proportional to its kinetic energy and the magnetic energy density:
$$t'_{cool} = \frac{6 \pi m_e c}{\sigma_T B'^2 \gamma}$$
The absolute maximum energy an electron can attain is managed by the precise equilibrium point where the acceleration timescale equals the synchrotron cooling timescale ($t'_{acc} \approx t'_{cool}$) (\citet{KumarandZhang2015}, p. 20).
Equating these two timescales isolates the maximum sustainable electron Lorentz factor:
$$\gamma_{max}^2 \approx \frac{6 \pi q}{\sigma_T B'}$$
One physical consequence emerges when this maximum electron energy is substituted into the standard synchrotron emission formula ($\nu'_{syn} \propto B' \gamma^2$). The magnetic field $B'$ cancels out, leaving a maximum comoving photon energy that relies just on fundamental constants. When this limiting comoving energy is translated into the observer frame, it establishes that even under perfect Bohm diffusion, the maximum synchrotron photon energy generated by shock accelerated electrons is restricted to:
$$E_{max} \approx 50 \Gamma \text{ MeV}$$
where $\Gamma$ is the bulk Lorentz factor of the relativistic outflow (\citet{KumarandZhang2015}, p. 20).
This $50 \Gamma \text{ MeV}$ limit serves as a critical diagnostic tool for astrophysicists: if photons are observed at energies much higher than this limit and no massive corresponding bulk Lorentz factor can be justified, the standard DSA framework in external shocks breaks down. Protons, because of their larger mass ($m_p/m_e \approx 1836$), suffer less from synchrotron cooling and can emit synchrotron photons up to $\sim 100 \Gamma \text{ GeV}$ before their acceleration is capped (\citet{KumarandZhang2015}, p. 20). However, resolving the maximum energy problem via hadronic acceleration introduces an efficiency crisis: because the total synchrotron power scales as $m^{-2}$, channeling the explosion's energy budget primarily into protons rather than electrons demands too high total kinetic energy to match the observed gamma ray luminosities.
If the cooling timescale is drastically shorter than the dynamical timescale of the expansion, electrons must be subjected to continuous acceleration deeper within the outflow to maintain the observed high energy continuum (\citet{KumarandZhang2015}, p. 105) and this requirement frequently shifts the theoretical preference away from localized shock fronts and toward macroscopic mechanisms, such as magnetic reconnection, which can maintain a hard electron distribution against catastrophic radiative losses (\citet{Peer2019}, p. 12).

\subsection{Ion temperatures and density profiles in extreme environments}
In conventional, lower temperature astrophysical plasmas, diagnostic spectroscopy relies much on the assumption of local thermodynamic equilibrium, utilizing Saha Boltzmann distributions to derive population level ratios and gas densities from atomic linewidth calculations. However, in the extreme environments of sub relativistic shocks and the innermost regions of advection dominated accretion flows, this classical equilibrium breaks down: electrons cool rapidly through intense synchrotron and inverse compton radiation, while heavier ions, experiencing significantly lower radiative energy losses, retain big thermal energies so that the ion temperature $T_i$ is different than the electron temperature, exceeding $10^{10} \text{ K}$ (\citet{Kafexhiuetal2019}, p. 1).
At these extreme thermal baselines, standard continuum fitting fails to constrain the physical state of the plasma and physicists must instead rely on discrete nuclear and hadronic signatures. The thermal broadening of prompt deexcitation lines, like the $4.438 \text{ MeV}$ emission from $^{12}\text{C}^*$, serves as a direct kinematic thermometer for the heavy ion population: when $T_i$ gets even higher, approaching the limit for inelastic proton-proton collisions, the spectrum becomes dominated by neutral pion decay. Since the $\pi^0 \rightarrow 2\gamma$ emissivity is strictly governed by the maxwellian velocity distribution of the constituent protons, the shape and overall amplitude of the resulting high energy gamma ray hump provide a measurement of the underlying proton temperature (\citet{Kafexhiuetal2019}, p. 1).
Quantifying the particle density of the emitting region is critical for understanding the outflow's mass loading and dynamic opacity; for localized, metal rich plasma blobs entrained within the relativistic jet, density extraction relies on ionization state mapping. As the intense radiation field from the central engine sweeps over these filaments, their opacity is dictated by the ionization parameter, $\xi = L / n r^2$, where $L$ is the ionizing luminosity, $n$ is the local particle density, and $r$ is the radial distance from the source (\citet{Haileyetal1999}, p. 1). By identifying specific K-shell absorption edges or blended resonant absorption troughs in the X-ray to soft gamma ray spectrum, observers can understand the prevailing value of $\xi$. If the emission radius and luminosity are constrained independently through temporal variability, this photoionization metric yields a precise calculation of the local density $n$ within the entrained clumps (\citet{Haileyetal1999}, p. 2).
On macroscopic scales, the density profile of the extended circumburst medium dictates the overall kinematic deceleration of the outflow. As the jet propagates outward, the efficiency with which it dissipates bulk kinetic energy through collective plasma instabilities, such as the two stream and filamentation instabilities, depends heavily on the ambient particle concentration. Theoretical models of these collisionless interactions show that the interstellar medium density must satisfy a firm lower limit to trigger gamma ray production, a threshold that scales with the bulk Lorentz factor $\Gamma$ \citep{Brainerd1999}. If a burst detonates in a highly rarefied environment that falls below this density floor, the instabilities fail to efficiently convert kinetic energy into thermalized electrons and these events undergo a structural selection effect, suppressing prompt gamma ray emission and manifesting only as lower energy ultraviolet or optical transients \citep{Brainerd1999}.
Furthermore, the large scale external density structure governs the onset and decay rate of the multi wavelength afterglow: the transition from the initial coasting phase to the deceleration phase occurs when the jet has swept up a volume of external mass more or less equivalent to its initial rest mass divided by $\Gamma_0$. The exact timing of this deceleration peak in the early light curve ($t_{dec}$) directly probes the radial density profile of the surrounding environment (\citet{KumarandZhang2015}, p. 30); researchers can distinguish between a uniform interstellar medium where $n$ remains constant and a stratified stellar wind profile where $n \propto r^{-2}$, characteristic of a massive Wolf-Rayet progenitor (\citet{KumarandZhang2015}, p. 32), by modeling how the characteristic synchrotron cooling frequencies evolve downstream of the forward shock. The observed high energy emissions thus function as continuous density traces, tracing the mass distribution from the immediate, highly ionized vicinity of the black hole out to the boundaries of the pre explosion stellar environment.

\section{Conclusions} 
Understanding how gamma ray emitting plasmas work requires a mix of microscopic particle interactions and macroscopic magnetohydrodynamic outflows. The observed spectra from things like gamma ray bursts and AGNs are not just endpoints of radiation, but they are records of energy partition: as seen throughout this review, the broad non thermal continua map the efficiency of shock acceleration and magnetic reconnection, while thermal and hadronic features map direct kinematic thermometers for extreme baryonic environments. The theoretical shift from kinetic Weibel instabilities toward Poynting flux dominated flows highlights the necessity of macroscopic magnetic fields in sustaining high energy particle distributions against catastrophic radiative cooling.
Meanwhile, the transition from isolated photon counting to a multi messenger framework has altered our view: joining orbital scale gravitational wave measurements and TeV neutrino detections with broadband electromagnetic data eliminates the physical degeneracies inherent in highly attenuated, distant sources. This allows astrophysicists to isolate the baryonic fraction of relativistic jets and constrain the bulk Lorentz factors required to overcome localized pair production opacity.
Furthermore, resolving the maximum energy problem and the 'synchrotron line of death' will require advanced spectroscopy to decouple quasi thermal photospheric emissions from non thermal shock signatures. The monitoring of transient decay slopes, chemical spallation markers, and propagation delays will refine our mapping of the circumburst medium. 


\bibliographystyle{plainnat} 
\bibliography{references}

\clearpage

\end{document}